% TOP_PROBLEM: {"id": 3132, "title": "Minimum number of arc-disjoint transitive subtournaments of order 3 in a tournament", "category_id": 3, "category": {"id": 3, "name": "graph_theory", "display_name": "Graph Theory", "description": "Problems involving graphs, networks, and their properties.", "slug": "graph-theory", "order_index": 3, "created_at": "2026-07-31T15:26:25.671Z"}, "status": "open", "problem_number": "OPG-49795", "set_id": 12, "statusQualification": "The auxiliary extremal function is the minimum, over tournaments, of the maximum packing size; the source discussion omits this inner maximum."}
% FIRST_POSTED: 2026-10-01
% ENTRY_KIND: statement_only
% UnsolvedMath ID 3132; source catalogue code OPG-49795
% !TeX program = lualatex
% Definition and sources only; this file is not an LLM solution attempt.
\documentclass[11pt,a4paper]{article}
\usepackage[margin=25mm]{geometry}
\usepackage{fontspec}
\setmainfont{lmroman10-regular.otf}[BoldFont=lmroman10-bold.otf,
 ItalicFont=lmroman10-italic.otf,BoldItalicFont=lmroman10-bolditalic.otf]
\usepackage{amsmath,amssymb,mathtools}
\usepackage{unicode-math}
\setmathfont{latinmodern-math.otf}
\AtBeginDocument{\let\setminus\smallsetminus}
\usepackage{xurl,hyperref}
\hypersetup{colorlinks=true,urlcolor=blue}
\setcounter{secnumdepth}{0}
\setlength{\parindent}{0pt}
\setlength{\parskip}{5pt}
\setlength{\emergencystretch}{3em}
\begin{document}
\section*{Minimum number of arc-disjoint transitive subtournaments of order 3 in a tournament}\label{3132}
{\small UnsolvedMath ID 3132 · OPG-49795 · Graph Theory · OpenGarden}
\subsection{Definitions and mathematical statement}
A tournament of order $n\geq1$ is a directed graph with $n$ vertices,
no loops, and exactly one of the arcs $u\to v$ and $v\to u$ for each
pair of distinct vertices. A transitive triple is a three-vertex
subtournament whose vertices can be ordered $a,b,c$ with arcs
$a\to b$, $a\to c$, and $b\to c$. Thus its three arcs do not form
a directed cycle.

A packing of transitive triples is a collection of such triples whose
arc sets are pairwise disjoint. Different triples may share a vertex, but no pair can share
two vertices, but cannot reuse any directed edge. Write $\nu_T(T)$
for the largest number of triples in a packing in $T$ and define
\[
 f(n)=\min\{\nu_T(T):T\text{ is a tournament on }n\text{ vertices}\}.
\]
The maximum inside this minimum is essential: an arbitrary packing
could always be empty.

The conjectured guarantee is that every tournament $T$ of order $n$
satisfies
\[
 \nu_T(T)\geq
 \left\lceil\frac{n(n-1)}6-\frac n3\right\rceil
 =\left\lceil\frac{n(n-3)}6\right\rceil.
\]
The ceiling denotes the least integer greater than or equal to its
argument. The assertion asks for one simultaneous packing of this
size; counting all transitive triples without controlling their
overlap is insufficient. Since each selected triple uses three arcs,
the scale of the proposed bound is one third of all arcs in the
tournament, with a linear correction.
\subsection{Short English statement}
Can every tournament pack at least the stated number of transitive triples without reusing an arc?
\subsection{Sources}
\begin{itemize}
\item[S1] ULAM.AI and the UnsolvedMath Contributors, supplied problems.json, record 3132 (OPG-49795), OpenGarden collection. Catalogue identity and source transcription; CC BY 4.0.
\url{https://huggingface.co/datasets/ulamai/UnsolvedMath}.
\item[S2] Open Problem Garden, Minimum number of arc-disjoint transitive subtournaments of order 3 in a tournament. Original problem and discussion; the mathematical formulation below is expanded from the supplied source record.
\url{http://www.openproblemgarden.org/op/minimum_number_of_transitive_subtournaments_of_order_3_in_a_tournament}.
\end{itemize}
\end{document}
